\documentclass[10pt,a4paper]{amsart}
\usepackage[margin=19mm]{geometry}
\usepackage{amsmath,amssymb,mathtools}
\usepackage[scr=rsfs,cal=euler]{mathalfa}
\usepackage{xfrac}
\usepackage[unicode,hidelinks,bookmarks=false]{hyperref}
\newcommand{\ie}{i.e.\ }
\newcommand{\cf}{cf.\ }
\newcommand{\resp}{respectively\ }
\newcommand{\Subsection}{\subsection}
\providecommand{\nobreakdash}{\nobreak\hbox{-}}
\setlength{\emergencystretch}{2em}
\multlinegap0pt
\newtheorem{prop}{Proposition}
\newcommand{\cc}{\circ}
\newcommand{\coz}{{\rm Coz}}
\newcommand{\de}{\delta}
\newcommand{\ov}{\overline}
\newcommand{\si}{\sigma}
\title{An extension of $F$-spaces and its applications}
\author{A.R. Aliabad, A. Taherifar}
\date{Source: arXiv:2509.07830v1 (CC BY 4.0)}
\begin{document}
\maketitle

\noindent\emph{Source and licence.} An extension of $F$-spaces and its applications. Version arXiv:2509.07830v1 (CC BY 4.0), distributed by arXiv under the Creative Commons Attribution 4.0 International licence, http://creativecommons.org/licenses/by/4.0/, as recorded in the arXiv licence metadata for this exact version and shown on the abstract page https://arxiv.org/abs/2509.07830v1. Full licence notice: \texttt{licenses/arxiv-2509.07830-CC-BY-4.0.txt}.\par\smallskip

\noindent\textit{Editorial scope.} Counted unit: the article's prop jan (source0.tex lines 403--406). The article's complete original proof of it is delivered with it (source0.tex lines 407--417, one \texttt{proof} environment of 11 lines). No mathematics is written, repaired or re-derived: the statement and the proof are the article's own lines, in the article's own notation, and no retained line is substituted or re-flowed.  No further source-local context is needed: every cross-reference of every retained line resolves inside the two delivered spans.  Nothing was omitted from any retained span, so the package carries no omission record.  No retained line contains a citation, so the wrapper carries no bibliography and no bibliography entry.  No retained line contains a figure environment, an image-inclusion command or drawing code, so no image is delivered and the package declares no asset dependency.  Editorial formatting only, in addition: an AMS \texttt{amsart} preamble that restates the article's own declarations and macros used by the retained lines, namely the environments \texttt{prop} on separate counters, together with the standard packages those lines need.  The retained environments print in this document as Proposition 1 in order of appearance, whereas the article prints the same items with the numbering of its own full document.  The complete native source of the article as distributed in the arXiv e-print for this exact version is bundled unchanged as \texttt{sources/184778/source0.tex}.
\begin{prop}\label{jan}
Let $X$ be a topological space with only one non-isolated point $\si$, where $\si$ is not a $G_\de$-point.  
Then $X$ is a $WCF$-space if and only if, for any two disjoint cozero-sets, one of them is a clopen subset.
\end{prop}

\begin{proof}
($\Rightarrow$)  
Suppose $A,B\in\coz(X)$ are disjoint. It suffices to show that $\si\notin \ov{A}\cap \ov{B}$.  
Assume, to the contrary, that $\si\in \ov{A}\cap \ov{B}$.  
Let $Z_1,Z_2\in Z(X)$ be such that $A\subseteq Z_1^\cc$ and $B\subseteq Z_2^\cc$.  
Thus $\si\in Z_1\cap Z_2 \in Z(X)$.  Since $\si$ is not a $G_\de$-point, the set $Z_1\cap Z_2$ must contain an isolated point.  
Hence $(Z_1\cap Z_2)^\cc\neq\emptyset$, which gives a contradiction.  

($\Leftarrow$)  
This direction is immediate.
\end{proof}

\end{document}
